\documentclass[11pt]{article}
\usepackage[margin=1in]{geometry}
\usepackage{amsmath,amssymb,amsthm,mathtools}
\usepackage{graphicx}
\usepackage{microtype}
\usepackage[colorlinks=true,linkcolor=blue,urlcolor=blue,citecolor=blue]{hyperref}

\newtheorem{theorem}{Theorem}
\newtheorem{proposition}[theorem]{Proposition}
\newtheorem{lemma}[theorem]{Lemma}
\newcommand{\Hh}{\mathcal H}
\newcommand{\Bcal}{\mathcal B}
\newcommand{\CP}{\mathsf{CP}}
\newcommand{\tr}{\operatorname{tr}}
\newcommand{\id}{\mathbf 1}
\newcommand{\Cw}{\mathsf C}
\newcommand{\vecg}{\operatorname{vec}}

\title{Separability is unnecessary for the unique decomposition of bounded CP-semigroup generators}
\author{Candidate full answer to the first open question in arXiv:2401.14344}
\date{27 August 2026}

\begin{document}
\maketitle

\noindent\fbox{\parbox{0.94\linewidth}{\textbf{Status: candidate full
solution, likely valid.}  The main bijectivity theorem of vom Ende remains
true on an arbitrary complex Hilbert space.  The proof uses the source's
weighted Choi kernel lemma one coordinate at a time and treats arbitrary
Kraus families as nets over finite subsets.}}

\section{Source question}

Vom Ende proves a unique decomposition theorem for generators of
norm-continuous completely positive semigroups on a \emph{separable} Hilbert
space \cite[Theorem~5.8 in the arXiv HTML; label \texttt{thm\_main} in the
source]{vomEnde}.  Section~6, PDF page~19, asks:
\begin{quote}
``An obvious question is whether separability of the underlying Hilbert space
is a necessary assumption in Theorem~5.1.''
\end{quote}
The PDF's production numbering displays the theorem and proposition as 5.1;
the arXiv HTML identifies them as Theorem~5.8 and Proposition~5.6.

\begin{figure}[ht]
\centering
\includegraphics[width=\linewidth]{figures/open_problem_crop.png}
\caption{The first open question in Section~6 of arXiv:2401.14344v4,
PDF page~19.}
\end{figure}

\section{Definitions and source tools}

Let $\Hh$ be a complex Hilbert space, with no separability assumption.
Write $\Bcal^1(\Hh)$ for the trace class and $\Bcal(\Hh)$ for the bounded
operators.  The cone $\CP(\Hh)$ consists of bounded completely positive maps
on $\Bcal^1(\Hh)$.  For $B\in\Bcal^1(\Hh)$ the source defines
\[
 \CP_B(\Hh)=\left\{\Phi\in\CP(\Hh):
 \tr\bigl(\Phi(B^*(\cdot)B)\bigr)=0\right\},
\]
where the displayed trace is the trace of the induced trace-class operator on
the Hilbert--Schmidt class.  By Lemma~4.7 of the source, if
\[
 \Phi(X)=\sum_{j\in J}V_jXV_j^*,
 \qquad \sup_{F\Subset J}\left\|\sum_{j\in F}V_j^*V_j\right\|<\infty,
\]
then $\Phi\in\CP_B(\Hh)$ exactly when
$\tr(B^*V_j)=0$ for every $j\in J$.  Here and below $F\Subset J$ means that
$F$ is finite, and sums over arbitrary $J$ are finite-subset nets.

We use one consequence of Lemmas~5.2 and~5.4 of the source.  Let
$G=\{g_j\}_{j\in J}$ be an orthonormal basis, let
$\lambda\in\ell^2(J)$, and put
\[
 D_\lambda=\sum_{j\in J}\lambda_j|g_j\rangle\langle g_j|.
\]
For a bounded map $\Psi$ let $\Cw_{\lambda,G}(\Psi)$ be the source's weighted
Choi operator.  If $X\in\Bcal^2(\Hh)$ and
$\Psi\in\CP_{XD_\lambda}(\Hh)$, then
\begin{equation}\label{eq:kernel-tool}
 \Cw_{\lambda,G}(\Psi)\vecg_G(X)=0.
\end{equation}
Crucially, this kernel statement is proved in the source for arbitrary
Hilbert spaces; it does not require that the weighted Choi map be injective.

\section{Proof intuition}

The separable proof chooses one square-summable diagonal weight which is
nonzero on every basis vector.  Such a weight cannot exist on an uncountable
basis.  But the proof only divides by the weight at the column currently
being inspected.  We therefore choose a new weight for each column: it is
supported on the countable singular support of $B$ together with that one
basis vector.  The source's kernel computation then determines that column of
$K$.  Repetition over the basis proves that $K$ is scalar.

For existence, a CP map on a nonseparable space may have an uncountable Kraus
index set.  The centering coefficients nevertheless form an $\ell^2(J)$
family.  Cauchy--Schwarz therefore makes the relevant finite-subset net
converge in operator norm and keeps the centered Kraus square sums uniformly
bounded.  This is the net version of the source's countable calculation.

\section{The uniqueness obstruction without separability}

\begin{proposition}\label{prop:obstruction}
Let $\Hh$ be an arbitrary complex Hilbert space,
$B\in\Bcal^1(\Hh)$ with $\tr(B)\ne0$, and $K\in\Bcal(\Hh)$.  If
\[
 K(\cdot)+(\cdot)K^*\in
 \CP_B(\Hh)-\CP_B(\Hh),
\]
then $K=i\alpha\id$ for some $\alpha\in\mathbb R$.
\end{proposition}

\begin{proof}
Choose a Schmidt decomposition
\[
 B=\sum_{n\in N}s_n|f_n\rangle\langle g_n|,
 \qquad s_n>0,
\]
where $N$ is countable because $B$ is trace class.  Extend
$\{g_n:n\in N\}$ to an orthonormal basis
$G=\{g_j\}_{j\in J}$ of $\Hh$.  Define the Hilbert--Schmidt operator
\[
 B_1=\sum_{n\in N}\sqrt{s_n}|f_n\rangle\langle g_n|.
\]

Fix one $k\in J$.  Define a nonnegative weight
$\lambda^{(k)}\in\ell^2(J)$ by
\[
 \lambda^{(k)}_n=\sqrt{s_n}\quad(n\in N),\qquad
 \lambda^{(k)}_k=1\quad(k\notin N),\qquad
 \lambda^{(k)}_j=0\quad(j\notin N\cup\{k\}).
\]
The clauses are disjoint, and
$\|\lambda^{(k)}\|_2^2\leq\|B\|_1+1$.  With
\[
 B_2^{(k)}=\sum_{j\in J}\lambda^{(k)}_j
 |g_j\rangle\langle g_j|
\]
we have $B_2^{(k)}\in\Bcal^2(\Hh)$,
$\lambda^{(k)}_k>0$, and, importantly,
\begin{equation}\label{eq:factorization}
 B_1B_2^{(k)}=B.
\end{equation}

Write
\[
 K(\cdot)+(\cdot)K^*=\Phi_1-\Phi_2,
 \qquad \Phi_1,\Phi_2\in\CP_B(\Hh).
\]
By \eqref{eq:factorization}, the kernel relation
\eqref{eq:kernel-tool} applies to each $\Phi_r$ with $X=B_1$ and
$D_\lambda=B_2^{(k)}$.  Hence it also applies to their difference.  Taking
the matrix coefficient at $g_k\otimes g_j$ and expanding the weighted Choi
operator exactly as in the proof of Proposition~5.6 of the source gives, for
every $j\in J$,
\begin{align}
0&=\left\langle g_k\otimes g_j,
 \Cw_{\lambda^{(k)},G}\bigl(K(\cdot)+(\cdot)K^*\bigr)
 \vecg_G(B_1)\right\rangle \notag\\
&=\lambda^{(k)}_k\left(
 \langle g_j,Kg_k\rangle\sum_{b\in J}\langle g_b,Bg_b\rangle
 +\delta_{jk}\sum_{b\in J}\langle g_b,K^*Bg_b\rangle
 \right)\notag\\
&=\lambda^{(k)}_k\left(
 \langle g_j,Kg_k\rangle\tr(B)
 +\delta_{jk}\tr(BK^*)\right).
\label{eq:column}
\end{align}
The sums are legitimate because $B$ and $BK^*$ are trace class.  Notice that
only $\lambda^{(k)}_k$, not weights at all basis coordinates, is divided out.

Since $\lambda^{(k)}_k>0$ and $\tr(B)\ne0$, equation
\eqref{eq:column} gives
\[
 \langle g_j,Kg_k\rangle=0\quad(j\ne k),\qquad
 \langle g_k,Kg_k\rangle
 =-\frac{\tr(BK^*)}{\tr(B)}=:c.
\]
The scalar $c$ is independent of $k$.  Since $k\in J$ was arbitrary,
$Kg_k=cg_k$ for every basis vector, hence $K=c\id$.

Finally, both $\Phi_1$ and $\Phi_2$ have zero $B$-weighted superoperator
trace.  Lemma~4.5(iv) of the source therefore gives
\[
0=\tr(KB^*)\tr(B)+\tr(B^*)\tr(BK^*)
  =2|\tr(B)|^2\operatorname{Re}c.
\]
Thus $\operatorname{Re}c=0$, proving $K=i\alpha\id$ with
$\alpha\in\mathbb R$.
\end{proof}

\section{Full nonseparable extension}

Let $\mathsf L(\CP(\Hh))$ denote the generators of norm-continuous
one-parameter semigroups of completely positive maps on
$\Bcal^1(\Hh)$.  The standard bounded-generator representation used in the
source says that every $L$ in this Lie wedge has the form
\begin{equation}\label{eq:standard}
 L=K_0(\cdot)+(\cdot)K_0^*+\Phi_0,
 \qquad K_0\in\Bcal(\Hh),\quad\Phi_0\in\CP(\Hh),
\end{equation}
and is valid without a separability hypothesis
\cite{ChristensenEvans,vomEnde}.

\begin{theorem}\label{thm:main}
Let $\Hh$ be an arbitrary complex Hilbert space and let
$B\in\Bcal^1(\Hh)$ satisfy $\operatorname{Re}\tr(B)\ne0$.  Then
\[
\begin{aligned}
\Xi_B:\{K\in\Bcal(\Hh):\operatorname{Im}\tr(B^*K)=0\}
 \times\CP_B(\Hh)&\longrightarrow\mathsf L(\CP(\Hh)),\\
(K,\Phi)&\longmapsto K(\cdot)+(\cdot)K^*+\Phi
\end{aligned}
\]
is bijective.  In particular, separability is not necessary in the source's
main theorem.
\end{theorem}

\begin{proof}
We first prove surjectivity.  Take $L$ and a representation
\eqref{eq:standard}.  The arbitrary-Hilbert-space Kraus theorem recalled in
Section~4 of the source gives an index set $J$ and operators
$\{V_j\}_{j\in J}\subset\Bcal(\Hh)$ such that
\[
 \Phi_0(X)=\sum_{j\in J}V_jXV_j^*,\qquad
 M:=\left\|\sum_{j\in J}V_j^*V_j\right\|<\infty.
\]
Put
\[
 a_j=\frac{\tr(B^*V_j)}{\tr(B^*)}.
\]
The denominator is nonzero, and Lemma~4.7 of the source gives
$a=(a_j)_{j\in J}\in\ell^2(J)$.  For $F\Subset J$ and $x\in\Hh$,
\[
 \left\|\sum_{j\in F}\overline{a_j}V_jx\right\|
 \leq\left(\sum_{j\in F}|a_j|^2\right)^{1/2}
      \left(\sum_{j\in F}\|V_jx\|^2\right)^{1/2}
 \leq\|a|_F\|_2\sqrt M\,\|x\|.
\]
Applied to tails, this shows that the finite-subset net converges in operator
norm.  Denote its limit and the squared norm of $a$ by
\[
 W=\sum_{j\in J}\overline{a_j}V_j,
 \qquad q=\sum_{j\in J}|a_j|^2.
\]

Set $\widetilde V_j=V_j-a_j\id$.  For every finite $F$ and $x\in\Hh$,
\[
 \sum_{j\in F}\|\widetilde V_jx\|^2
 \leq2\sum_{j\in F}\|V_jx\|^2
     +2\sum_{j\in F}|a_j|^2\|x\|^2
 \leq2(M+q)\|x\|^2.
\]
Thus the centered family defines a completely positive map
\[
 \Phi(X)=\sum_{j\in J}\widetilde V_jX\widetilde V_j^*.
\]
Moreover,
$\tr(B^*\widetilde V_j)=\tr(B^*V_j)-a_j\tr(B^*)=0$,
so $\Phi\in\CP_B(\Hh)$ by Lemma~4.7 of the source.

Expanding finite sums and passing to the trace-norm limit gives
\[
 \Phi(X)=\Phi_0(X)-WX-XW^*+qX
 \qquad(X\in\Bcal^1(\Hh)).
\]
Indeed, the Kraus sums converge in trace norm by Lemma~4.3 of the source,
$W_F\to W$ in operator norm, and the scalar partial sums tend to $q$.
Consequently, with
\[
 \widetilde K=K_0+W-\frac q2\id,
\]
we have
\[
 L=\widetilde K(\cdot)+(\cdot)\widetilde K^*+\Phi.
\]
Finally put
\[
 K=\widetilde K-i\,
 \frac{\operatorname{Im}\tr(B^*\widetilde K)}
      {\operatorname{Re}\tr(B)}\,\id.
\]
The added imaginary scalar does not change
$K(\cdot)+(\cdot)K^*$, and direct calculation gives
$\operatorname{Im}\tr(B^*K)=0$.  This proves surjectivity.

For injectivity, suppose
$\Xi_B(K_1,\Phi_1)=\Xi_B(K_2,\Phi_2)$.  Then
\[
 (K_2-K_1)(\cdot)+(\cdot)(K_2-K_1)^*=\Phi_1-\Phi_2.
\]
Proposition~\ref{prop:obstruction} gives
$K_2-K_1=i\alpha\id$ for some real $\alpha$.  The two domain gauge
conditions imply
\[
0=\operatorname{Im}\tr\bigl(B^*(K_2-K_1)\bigr)
 =\alpha\operatorname{Re}\tr(B),
\]
so $\alpha=0$.  Hence $K_1=K_2$, and the displayed generator identity then
gives $\Phi_1=\Phi_2$.
\end{proof}

The source's corollaries for bounded quantum-dynamical semigroup generators
are obtained formally from the main theorem and therefore also lose their
separability hypothesis.

\section{Verification, scope, and novelty}

The proof uses only results which the source states for arbitrary Hilbert
spaces: the weighted Choi kernel identity, the arbitrary-index Kraus
representation and convergence lemma, and the trace formula for centered
Kraus operators.  The new point in Proposition~\ref{prop:obstruction} is that
the Choi weight varies with $k$; no injectivity of the weighted Choi map is
claimed or needed.  The existence proof explicitly uses nets indexed by
finite subsets, so it does not hide a separability assumption.

No computational evidence is used.  The main expert-review targets are
equation~\eqref{eq:column}, the passage from arbitrary-index Kraus sums to the
centered map, and the separability-free scope of the standard representation
\eqref{eq:standard}.

Bounded searches on 27 August 2026 covered the run's four lightweight indexes,
arXiv/web searches for the exact question, the paper title and author,
``Proposition 5.6'', ``nonseparable'', and close CP-generator-decomposition
phrases.  They found the source, its publication record, and the author's
finite-dimensional predecessor, but no later solution or equivalent
nonseparable extension.  This supports only moderate novelty confidence: the
search was not exhaustive, and the coordinatewise-weight observation is
elementary once isolated.

\begin{thebibliography}{9}
\bibitem{vomEnde}
F.~vom Ende,
\emph{From the Choi Formalism in Infinite Dimensions to Unique
Decompositions of Generators of Completely Positive Dynamical Semigroups},
arXiv:2401.14344v4 (2024),
\url{https://arxiv.org/abs/2401.14344}.

\bibitem{ChristensenEvans}
E.~Christensen and D.~E.~Evans,
\emph{Cohomology of operator algebras and quantum dynamical semigroups},
J. London Math. Soc. (2) \textbf{20} (1979), 358--368.
\end{thebibliography}

\end{document}

